\section{Section 1: Foundational Architectural Principles: A Co-Designed Approach for Chaotic Systems}

The accurate and long-term simulation of $\mathrm{N}$-body gravitational systems, particularly the notoriously chaotic three-body problem, presents a formidable challenge for traditional numerical and machine learning methods alike. The dynamics are characterized by an extreme sensitivity to initial conditions, where infinitesimal perturbations can lead to exponentially divergent trajectories over time. Furthermore, the solutions often contain a complex mixture of smooth, quasi-periodic motion and sharp, high-frequency events that occur during close encounters between bodies. A successful predictive model must therefore not only capture these multi-scale temporal features with high fidelity but also adhere to fundamental global conservation laws, such as the conservation of total energy and angular momentum, to ensure its long-term physical plausibility.$^1$

A generic neural network architecture, such as a standard Multi-Layer Perceptron (MLP), lacks the specific structural priors necessary to handle this confluence of challenges. Such models often suffer from a "spectral bias," a tendency to learn low-frequency functions more easily, which makes them struggle to represent the high-frequency components inherent in chaotic dynamics.$^2$ More critically, enforcing the governing physical laws—in this case, Newton's laws of motion—within a Physics-Informed Neural Network (PINN) framework requires the computation of high-order derivatives of the network's output. For the $\mathrm{N}$-body problem, this involves calculating the second derivative of the predicted position with respect to time, $\frac{\mathrm{d}^2 \mathrm{r}}{\mathrm{d} t^2}$, to form the physics residual term in the loss function. Common activation functions like the Rectified Linear Unit (ReLU) are fundamentally unsuitable for this task. While computationally efficient, ReLU produces a piecewise linear output, meaning its first derivative is piecewise constant and its second derivative is a series of Dirac delta functions (i.e., zero almost everywhere). This mathematical property makes it impossible to represent the smooth, continuously varying gravitational forces that govern celestial motion.$^1$

To overcome these fundamental limitations, a more principled approach is required—one where the architecture itself imparts a strong inductive bias that aligns with the underlying physics. The framework detailed in this report proposes a hybrid PINN that is not an arbitrary collection of techniques but a meticulously co-designed system. Each architectural component is deliberately chosen to counteract a specific, well-understood failure mode of standard models when applied to chaotic physical systems. This synergistic design methodically addresses the core requirements of the problem:

\begin{enumerate}
    \item \textbf{Derivative Fidelity:} To accurately enforce the physics residual, the network must be built with an activation function that is infinitely differentiable ($\mathrm{C}^\infty$) and whose derivatives are well-behaved. The Sinusoidal Representation Network (SIREN) architecture, which uses the sine function as its nonlinearity, is the ideal choice. The derivative of a sinusoid is another sinusoid, merely phase-shifted, a property that is preserved for higher-order derivatives. This provides a powerful inductive bias towards learning smooth functions whose derivatives are also smooth, a property essential for minimizing the physics residual with high fidelity.$^1$
    \item \textbf{Expressivity and Training Stability:} Capturing the intricate complexity of chaotic trajectories requires a network with significant expressive capacity, which is typically achieved through depth. However, training very deep neural networks is notoriously difficult due to the well-known problems of vanishing and exploding gradients. The Residual Network (ResNet) backbone provides a solution by incorporating "skip" connections that create an identity mapping shortcut, allowing gradients to flow unimpeded through the network's depth. This report specifies the use of a particularly effective variant, the Power-Enhanced SkipResNet (SQR-SkipResNet), which has been empirically shown to provide more stable gradient flow and superior performance on non-smooth functions, making it well-suited for deep architectures.$^1$
    \item \textbf{Long-Term Physical Plausibility:} Even a PINN that perfectly satisfies the local, instantaneous equations of motion can accumulate numerical drift over long time horizons, leading to unphysical outcomes such as a gradual increase or decrease in the total energy of an isolated system. While specialized architectures like Hamiltonian Neural Networks (HNNs) can enforce energy conservation by construction through a symplectic structure, they often impose rigid constraints on the network's form (e.g., requiring the use of canonical coordinates) that can limit its flexibility as a general function approximator.$^1$ The proposed architecture adopts a more pragmatic and flexible approach by incorporating the principle of energy conservation as a "soft" constraint within the loss function. This HNN-style energy regularizer penalizes the network for any deviation of the predicted total energy from its known initial value, guiding the optimization process toward solutions that lie on or near the constant-energy manifold defined by the initial conditions. This strategy provides the primary benefit of an HNN—energy conservation—without its restrictive architectural costs, thereby improving the long-term physical plausibility of the model's predictions.$^1$
\end{enumerate}

This philosophy of pragmatic hybridization, which strategically borrows the most critical benefits from multiple advanced paradigms while avoiding their most restrictive drawbacks, is a central theme of the proposed design. The resulting model is a synergistic whole, where the SQR-SkipResNet enables the necessary depth for expressivity, the SIREN activations provide the derivative fidelity required for physics enforcement, and the HNN-style loss term ensures the long-term stability essential for physical realism.
